% !TEX root = ../manuscript.tex

\section{Méthodologie et organisation}

\subsection{Approche CRISP-DM}

Ce projet a été conduit en suivant le cycle \textbf{CRISP-DM}~\cite{chapman2000}, qui structure un projet de data science en six phases~: compréhension du problème, compréhension des données, préparation des données, modélisation, évaluation, et déploiement. Chaque phase a alimenté la suivante de manière itérative~: l'EDA a révélé l'absence de signal linéaire, ce qui a orienté le choix des modèles non-linéaires~; la détection du data leakage a conduit à réviser le feature engineering~; les résultats de classification ont confirmé le diagnostic de l'EDA.

\subsection{Outils et environnement technique}

\begin{table}[h]
\centering
\caption{Stack technique du projet}
\label{tab:tech_stack}
\begin{tabular}{lll}
\toprule
Catégorie & Outil & Usage \\
\midrule
Langage              & Python 3                  & Ensemble du projet \\
Manipulation données & pandas, numpy             & Chargement, nettoyage, feature engineering \\
Modélisation         & scikit-learn, xgboost     & Entraînement, tuning, évaluation \\
Déséquilibre classes & imbalanced-learn          & Rééchantillonnage SMOTE \\
Interprétabilité     & shap                      & Valeurs SHAP sur XGBoost \\
Réduction de dim.    & scikit-learn (PCA)        & Visualisation clustering \\
Versionnement        & Git / GitHub              & Suivi des notebooks et du rapport \\
Rédaction            & \LaTeX                    & Mémoire \\
\bottomrule
\end{tabular}
\end{table}

\subsection{Compétences mobilisées}

\begin{table}[h]
\centering
\caption{Lien entre enseignements MIASHS et applications projet}
\label{tab:competences}
\begin{tabular}{ll}
\toprule
Enseignement M1 & Application dans le projet \\
\midrule
ACP / ACM                          & Réduction dimensionnelle pour la visualisation des clusters \\
Régression linéaire \& diagnostic  & Baseline régression, Ridge, Lasso, ElasticNet \\
Classification supervisée          & Logistic Regression, Random Forest, XGBoost \\
Apprentissage non supervisé        & K-Means, CAH \\
Descente de gradient \& optimisation & Compréhension de XGBoost et de la log-loss \\
Régression logistique              & Baseline classification + interprétation coefficients \\
Régularisation \& optimisation     & Ridge, Lasso, ElasticNet \\
Déséquilibre de classes            & class\_weight, SMOTE (imbalanced-learn) \\
Python                             & Ensemble du pipeline de data science \\
\bottomrule
\end{tabular}
\end{table}

\section{Difficultés rencontrées et retour d'expérience}

\subsection{Difficultés techniques}

\textbf{La limite fondamentale du dataset synthétique} est la difficulté principale que j'ai rencontrée. Découvrir que $R^2 \approx 0$ et AUC $\approx 0.5$ pour tous les modèles a demandé un vrai travail de diagnostic~: relier ce résultat au constat de l'EDA (corrélations $< 0.01$) et à la nature synthétique des données a été l'exercice analytique central de ce projet.

\textbf{La détection du data leakage} sur \texttt{hypercholesterol}~\cite{kaufman2012} est une difficulté classique en ML, rarement rencontrée dans les travaux académiques. Le $R^2$ artificiel de 0.675 aurait pu passer inaperçu sans une analyse rigoureuse des feature importances.

\textbf{Le coût computationnel}~: un premier essai de Random Forest sans \texttt{max\_depth} ni \texttt{n\_jobs=-1} a conduit à un entraînement de plus de 3 heures interrompu avant convergence. Ça m'a forcé à comprendre les paramètres internes de l'algorithme. Le SVM a été écarté pour la même raison, sa complexité en $\mathcal{O}(n^2)$ étant incompatible avec 100~000 observations.

\subsection{Compétences acquises}

\begin{itemize}
  \item \textbf{Détection et correction du data leakage}~\cite{kaufman2012}~: piège courant non systématiquement couvert en cours théoriques.
  \item \textbf{RandomizedSearchCV}~\cite{bergstra2012}~: logique de l'exploration aléatoire vs exhaustive des hyperparamètres.
  \item \textbf{Gestion du déséquilibre de classes}~: comparaison concrète entre \texttt{class\_weight='balanced'} et SMOTE~\cite{chawla2002}, et compréhension de la différence entre accuracy, F1-score et AUC-ROC~\cite{sokolova2009}.
  \item \textbf{Interprétabilité multi-méthodes}~: la divergence entre outils est elle-même informative~\cite{rudin2019}.
  \item \textbf{Analyse critique d'un résultat négatif}~: savoir diagnostiquer pourquoi un modèle échoue est aussi précieux que d'en construire un qui réussit.
\end{itemize}

\subsection{Plus-value~: diagnostiquer les biais du pipeline}

Ce projet ne produit pas un modèle prédictif déployable --- les métriques ($R^2 \approx 0$, AUC $\approx 0.5$) l'excluent. Sa plus-value réside dans une contribution différente~: \textbf{construire un pipeline capable de diagnostiquer pourquoi un modèle échoue}, et d'identifier rigoureusement ses biais.

Trois apports concrets peuvent être retenus~:

\begin{enumerate}
  \item \textbf{Détection du data leakage.} La variable \texttt{hypercholesterol} produisait un $R^2$ artificiel de 0.675~\cite{kaufman2012}. Un pipeline mal construit aurait publié ce résultat comme une performance réelle. La détection et la correction rigoureuse de ce biais est une compétence directement valorisable sur des données épidémiologiques réelles.

  \item \textbf{Diagnostic de l'absence de signal par cinq preuves convergentes.} Corrélations de Pearson $< 0.01$ (EDA), $R^2 \approx 0$ sur six modèles de régression dont des méthodes non-linéaires (Random Forest, XGBoost), AUC $\approx 0.5$ sur trois modèles de classification, silhouettes $\approx 0.10$ en clustering, et divergence des cinq outils d'interprétabilité. Cette convergence de preuves indépendantes est plus robuste qu'un résultat isolé.

  \item \textbf{Accuracy trompeuse sur données déséquilibrées.} Un modèle naïf atteignait 75\% d'accuracy sur le déséquilibre 75/25~; le Random Forest affichait 56.8\% d'accuracy mais une AUC de 0.494, soit \textit{pire} que le hasard pur~\cite{sokolova2009}. Identifier et exposer ce piège méthodologique est un résultat analytique à part entière.
\end{enumerate}

Ce pipeline est \textbf{directement transposable} sur des données réelles~\cite{dawber1951, nhanes}~: les mêmes étapes de détection de leakage, de validation des métriques et de comparaison des méthodes s'appliquent, et le signal réel permettrait d'obtenir des résultats prédictifs exploitables.

\subsection{Biais identifiés et enjeux éthiques}

Au-delà du data leakage --- biais technique identifié et corrigé --- ce projet met en évidence des biais plus fondamentaux, de nature méthodologique et éthique, qu'il est nécessaire d'analyser explicitement.

\textbf{1. Biais de représentation~: l'impossibilité de généraliser.} Le dataset est synthétique~: les 100~000 individus ne représentent aucune population réelle. Les distributions ont été générées sans ancrage épidémiologique validé (aucune cohorte de référence, aucun rapport sur la méthode de génération). Cela rend toute généralisation structurellement impossible~: un modèle entraîné ici ne peut pas être évalué sur Framingham~\cite{dawber1951} ou NHANES~\cite{nhanes}, car les distributions diffèrent de façon inconnue. Ce biais est antérieur à toute modélisation~: il est inscrit dans le choix du dataset.

\textbf{2. Biais de conception~: répondre à une question avec des données inadaptées.} L'objectif initial est de modéliser les relations entre habitudes de vie et indicateurs biologiques. Or, dans un dataset synthétique sans structure causale, les variables ont été générées indépendamment~: il n'existe, par construction, aucune relation à découvrir~\cite{scholkopf2021}. Appliquer des méthodes de ML sur ces données revient à chercher un signal que l'on sait absent. Ce biais de conception est distinct du biais de représentation~: même si le dataset représentait une vraie population, l'absence de structure causale rendrait l'inférence impossible. En termes d'apprentissage causal~\cite{pearl2009}, le DAG sous-jacent est un graphe sans arcs --- chaque variable est une île.

\textbf{3. Enjeux de fairness algorithmique.} Si ce pipeline était déployé sur des données réelles, un biais critique émergerait~: les performances peuvent varier considérablement selon les sous-groupes. Le travail d'Obermeyer et al.~\cite{obermeyer2019} a démontré qu'un algorithme utilisé dans des millions d'hôpitaux américains présentait un biais racial systématique, pénalisant les patients noirs à sévérité médicale égale. En santé, la fairness n'est pas un aspect secondaire~: un modèle avec AUC globale acceptable peut masquer des performances très inégales selon le genre, l'âge, ou d'autres attributs protégés~\cite{rudin2019}. Dans ce projet, le déséquilibre 75/25 sur \texttt{disease\_risk} est lui-même un signal~: si cette variable était réelle, le seuil de décision du classificateur devrait être calibré en tenant compte de l'impact différentiel sur les groupes sous-représentés.

\textbf{4. Risques de déploiement~: quand l'absence de signal devient dangereuse.} Un résultat central de ce projet --- AUC $\approx 0.5$ --- signifie que les modèles ne font pas mieux que le hasard pour prédire \texttt{disease\_risk}. Ce résultat, s'il était mal interprété ou ignoré, pourrait conduire à des décisions erronées. Un praticien ou un décideur qui fait confiance à un modèle avec AUC = 0.494 prend des décisions au hasard, avec en plus l'illusion d'une base scientifique. Cette situation --- modèle inutile mais présenté comme fiable --- est l'un des risques les plus documentés de l'IA en santé~\cite{rajpurkar2022}. La calibration des prédictions (est-ce que "70\% de probabilité de risque" en sortie correspond vraiment à 70\% de cas réels~?) est une exigence minimale avant tout déploiement.

Ces quatre biais --- représentation, conception, fairness, déploiement --- constituent l'essentiel du travail critique qu'un data scientist doit conduire avant de présenter des résultats, y compris des résultats négatifs. Les identifier n'invalide pas le projet~: cela en précise le périmètre de validité et les conditions d'un travail de suivi crédible.

\subsection{Positionnement RSE / DDRS}

L'utilisation d'un dataset \textbf{synthétique} s'inscrit dans une démarche éthique~: les données de santé sont sensibles et encadrées par le RGPD~\cite{rgpd2016}. Les données synthétiques éliminent tout risque de ré-identification tout en permettant de développer des méthodes~--- un compromis central de l'IA responsable~\cite{rgpd2016}.

Le coût énergétique des entraînements a été limité en contraignant la profondeur des arbres (\texttt{max\_depth=15}) et en réduisant \texttt{n\_iter} dans \texttt{RandomizedSearchCV}, avec un impact positif sur l'empreinte carbone du calcul.

\subsection{Limites méthodologiques}

\textbf{Corrélations non-linéaires non explorées.} L'analyse de Pearson a révélé des coefficients inférieurs à $|0.01|$ --- mais Pearson ne détecte que les relations \textit{linéaires}~\cite{james2021}. Une analyse complémentaire aurait pu mobiliser le coefficient de Spearman (corrélation des rangs, capte les relations monotones), la \textit{mutual information} (\texttt{mutual\_info\_classif} dans scikit-learn, toute dépendance statistique quelle que soit sa forme), ou des nuages de points pour identifier visuellement des structures non-linéaires (relations en U, effets seuil, interactions).

Toutefois, l'utilisation de \textbf{Random Forest et XGBoost}~\cite{breiman2001} --- algorithmes capables de capturer n'importe quelle relation non-linéaire --- constitue une réponse indirecte à cette limite. Leur échec ($R^2 \approx 0$, AUC $\approx 0.5$) démontre que même les dépendances non-linéaires sont absentes~: résultat plus fort qu'un test de Spearman ou de mutual information, car il exploite toute la puissance des méthodes ensemblistes sur 100~000 individus.

\section{Perspectives}

\textbf{1. Changer de dataset}~: reproduire le pipeline sur un dataset réel, comme la cohorte Framingham~\cite{dawber1951} ou les données NHANES~\cite{nhanes}, qui contiennent de vraies relations entre habitudes de vie et indicateurs de santé. Les méthodes développées ici sont directement transposables.

\textbf{2. Améliorer les modèles de classification}~: explorer le \textbf{stacking} (combinaison de plusieurs modèles~\cite{hastie2009}) et les \textbf{réseaux de neurones} pour capturer des relations très non-linéaires~--- en lien direct avec le module Deep Learning du master. SMOTE ayant déjà été intégré dans ce projet, ces pistes constituent la suite logique.

\textbf{3. Déploiement}~: un modèle de classification pourrait être déployé sous forme d'API REST (FastAPI) permettant de soumettre un profil individuel et d'obtenir une prédiction avec ses valeurs SHAP associées, rendant le modèle accessible et explicable.

\textbf{4. Interprétabilité avancée}~: les \textit{Partial Dependence Plots} (PDP) et les \textit{Individual Conditional Expectation} (ICE) permettraient de visualiser comment la prédiction varie en fonction d'une variable, toutes choses égales par ailleurs~--- une limite de SHAP qui agrège les contributions.

\textbf{5. Robustesse sur le long terme (\textit{data drift})~:} sur des données réelles en production, un modèle se dégrade si la distribution des données évolue au fil du temps~--- on parle de \textit{concept drift}~\cite{hastie2009}. Les habitudes de vie d'une population changent, les seuils médicaux sont révisés, les profils épidémiologiques se transforment. En pratique, cela implique de surveiller les performances du modèle périodiquement (suivi de l'AUC sur de nouvelles cohortes) et de le réentraîner dès que les performances se dégradent. Cette problématique ne s'applique pas au dataset synthétique statique de ce projet, mais sera centrale sur des données réelles.

\textbf{6. Vers le Master 2}~: les thématiques de causalité, fairness algorithmique et apprentissage fédéré constituent des pistes naturelles pour un travail de recherche en M2, dans la continuité directe de ce projet.
